\begin{introduction}
    \item 函数
    \item 函数图像
    \item 函数的性质
    \item 反函数
    \item 一元二次函数
\end{introduction}

\begin{definition}[函数]
    函数是定义域和陪域都为数集的映射，例如：$f:\mathbb{R}\to \mathbb{R}$。

    由于都是数值，习惯把像称为因变量$y$，把原像称为自变量$x$，记作：$y=f(x)$。
\end{definition}

\section{函数的定义域}

\begin{definition}[定义域]
    函数的定义域是指自变量$x$的取值范围，即所有使得函数有意义的$x$的集合，记作$D_f$。
\end{definition}

\begin{theorem}[常用的定义域]
    求函数的定义域需要从这几个方面入手：

    \begin{minipage}[t]{0.65\linewidth}
        \begin{enumerate}
            \item 分母不为零。$\frac{1}{x-2} \Rightarrow x \neq 2$.
            \item 偶次根式的被开方数非负。$x^{\frac{1}{n}}=\sqrt[n]{x}, n \in 2k, \Rightarrow x \geqslant 0$.
            \item 对数中的真数部分大于 0. $\log_a x \Rightarrow x > 0$.
            \item 指数、对数的底数大于 0 ，且不等于 1。$\log_a x \Rightarrow a>0, a \neq 1$.
        \end{enumerate}
    \end{minipage}
    \hfill
    \begin{minipage}[t]{0.33\linewidth}
        \begin{enumerate}\setcounter{enumi}{4}
            \item $y=\tan x$ 中 $x \neq k \pi+\pi / 2$。
            \item $y=\cot x$ 中 $x \neq k \pi$ 等等。
            \item $x^0$ 中 $\mathrm{x} \neq 0$。
        \end{enumerate}
    \end{minipage}
\end{theorem}

\subsection{抽象函数的定义域}

\begin{example}
    设 $f(x+1)$ 的定义域为 $[0, a](a>0)$ ，则 $f(x)$ 的定义域为 \underline{\qquad}.
\end{example}

\v{4em}

\begin{example}
    $f(x)$ 定义域 $[0,2 a](a>0)$, 则 $F(x)=f(x)+f(x+a)$ 的定义域为 \underline{\qquad}.
\end{example}

\v{4em}

